\documentclass{article}
\usepackage[paper=a4paper,margin=1in]{geometry}

\usepackage[hidelinks]{hyperref}

\usepackage{xurl}

\usepackage[utf8]{inputenc}

\usepackage{subfiles}

\usepackage{minted}
\setminted{	
	numbersep=5pt,             % distanza numeri-codice
	frame=single,              % cornice
	breaklines=true,           % a capo automatico
	bgcolor=bggray	           % sfondo
}

\usepackage{xcolor}
\definecolor{bggray}{rgb}{0.97,0.97,0.97} % colore sfondo

\begin{document}


\begin{titlepage}
    \centering
    {\scshape\LARGE Gruppo 31\par}
    \vspace{1cm}
    {\scshape\Large PSS\par}
    \vspace{1.5cm}
    {\huge\bfseries Assignment 3 - Clean Code\par}
    \vspace{2cm}   
\end{titlepage}


\tableofcontents
\newpage

\section{Applicazzione scelta}
L'applicazione scelta per il progetto consiste in una piattaforma di prenotazione eventi TicketMaster, sviluppata in ambiente Android.

Questa permette ad un utente di creare un account, effetuare il login e cercare eventi in base a filtri e categorie; una volta selezionato un evento, l'utente potrà procedere
all'acquisto o prenotazione dei biglietti.

Per ottenere questi risultati l'applicazione si interfaccia con la API di TickerMaster, che fornisce i dati relativi agli eventi in formato JSON.

L'applicazione originale è disponibile su GitHub al seguente link: \url{https://github.com/FrancescoRomeo02/TicketWave_AndroidApp}.

    \subsection{Classi scelte}
    Le classi scelte per l'analisi sono le seguenti:
    \begin{itemize}
        \item \texttt{User.java}: rappresenta l'utente dell'applicazione, gestendo le sue proprietà e occupandosi dello scambio di informazioni legate a questo;
        \item \texttt{JsonParser.java}: si occupa della comunicazione con la API di TicketMaster, ottenendo i dati in formato JSON e parsandoli per estrarre le informazioni rilevanti;
        \item \texttt{EventAdapter.java}: gestisce la visualizzazione degli eventi all'interno della applicazione, crea viste dinamiche per mostrare gli eventi agli utenti;
        \item \texttt{DBHelperLiked.java}: ha il compito di gestire gli eventi preferiti dell'utente all'interno del database SQLite locale;
        \item \texttt{MainActivity.java}: rappresenta la classe principale, si occupa della procedura di login e navigazione tra le schermate dell'app.
    \end{itemize}

    \subsection{Gruppo di lavoro}
    Il gruppo è composto dai seguenti membri:
    \begin{itemize}
        \item Samuele Marelli - 899768
        \item Andrea Fumagalli - 900329
        \item Marco Passoni - 885873
    \end{itemize}

\newpage
\subfile{capitoli/EventAdapter.tex}
\subfile{capitoli/JsonParser.tex}
\subfile{capitoli/User.tex}

\end{document}